\documentclass[11pt]{article}
\usepackage{mathblatt}
\begin{document}
\blattkopf*{Basisaufgaben}{BAS-Z3}{Basisaufgaben · Zettel · BAS-Z3}
\einheitenkopf[][Zettel 3]{Basisaufgaben · Zettel 3}
\zweigzeile{ohne Taschenrechner · je Aufgabe 1 Punkt · Lösungen auf der Rückseite}
% \small: zehn Aufgaben mit bis zu zwei Koordinatensystemen auf einer Seite (Render 28.09.)
\small

% 1: Bruchteil einer Fläche bestimmen – brueche-dezimalzahlen-basis-k1-v3 (Original 2026-FOR-B1b)
\begin{aufgabe}{Bei welcher Figur ist genau $\frac{3}{8}$ grau? \hfill (P26F) \\ \kreuz{P} \kreuz{Q} \kreuz{R} \kreuz{S}}

P \bruchkreis[0.8]{3}{8} \quad Q \bruchrechteck[2.5]{3}{5} \quad R \kreissektor[0.8]{90}{} \quad S \bruchrechteck[2.5]{3}{10}
\end{aufgabe}

% 2: Zahlen in verschiedenen Darstellungen vergleichen – brueche-dezimalzahlen-basis-k4-v3 (Original 2023-OS-B1f)
\begin{aufgabe}{Welcher Wert ist der größte: $0{,}2$; $2\,\%$; $0{,}2^2$; $\frac{1}{4}$? \hfill (P23) \\ \leerfeld}
\end{aufgabe}

% 3: Termwert berechnen – rationale-zahlen-basis-k1-v2 (Original 2026-FOR-B1g)
\begin{aufgabe}{Berechne den Wert von $2 \cdot (x - 7)$ für $x = -1$. \hfill (P26F) \\ \leerfeld}
\end{aufgabe}

% 4: Prozentwert berechnen – prozentrechnung-basis-k4-v2 (Original 2026-FOR-B1a)
\begin{aufgabe}{$30\,\%$ von $60\,€$ \hfill (P26F) \\ \leerfeld[€]}
\end{aufgabe}

% 5: Lineare Gleichung lösen – lineare-gleichungen-basis-k1-v2 (Original 2024-OS-B1d)
\begin{aufgabe}{Löse die Gleichung $5 \cdot (x - 2) = 15$. \hfill (P24) \\ x = \leerfeld}
\end{aufgabe}

% 6: Graph einer linearen Funktion erkennen – lineare-funktionen-basis-k1-v1 (Original 2025-OS-B1f)
\begin{aufgabe}{\begin{minipage}[t]{0.6\linewidth}Welcher Graph gehört zu einer linearen Funktion? Kreuze an. \hfill (P25) \\ \kreuz{A} \kreuz{B} \kreuz{C} \kreuz{D}\end{minipage}\hfill\begin{minipage}[t]{0.37\linewidth}\vspace{0pt}
\begin{ksys}[xmin=-5,xmax=5,ymin=-5,ymax=5,klein] \parabel{1}{1}{-3}{A} \gerade[2.5]{0.5}{-2}{B} \funktionab[3.5]{3/\x}{C}{0.6}{5} \funktion{abs(\x)-3}{D} \end{ksys}
\end{minipage}}
\end{aufgabe}

% 7: Symmetrieachsen bestimmen – symmetrie-abbildungen-basis-k3-v1 (Original 2022-OS-B1i)
\begin{aufgabe}{\begin{minipage}[t]{0.6\linewidth}Eine Tischplatte ist ein Rechteck (kein Quadrat). Wie viele Symmetrieachsen hat die Figur? Kreuze an. \hfill (P22) \\ \kreuz{0} \kreuz{1} \kreuz{2} \kreuz{3} \kreuz{4} \kreuz{5}\end{minipage}\hfill\begin{minipage}[t]{0.37\linewidth}\vspace{0pt}
\rechteck{2.8}{1.6}
\end{minipage}}
\end{aufgabe}

% 8: Kreissektor Anteil berechnen – kreis-basis-k1-v1 (Original 2025-OS-B1e)
\begin{aufgabe}{\begin{minipage}[t]{0.6\linewidth}Im Bild ist ein Kreisausschnitt grau gefärbt. Der Mittelpunktswinkel des grauen Teils beträgt $90^\circ$. Wie viel Prozent der Kreisfläche sind grau? \hfill (P25) \\ \leerfeld[\%]\end{minipage}\hfill\begin{minipage}[t]{0.37\linewidth}\vspace{0pt}
\kreissektor[1.2]{90}{$90^\circ$}
\end{minipage}}
\end{aufgabe}

% 9: Kantenzahl eines Körpers angeben – koerper-basis-k2-v1 (Original 2019-OS-B1f)
\begin{aufgabe}{Ein Spielwürfel ist ein Würfel. Gib an, wie viele Kanten er hat. \hfill (P19) \\ \leerfeld Kanten}
\end{aufgabe}

% 10: Wahrscheinlichkeit einstufig – bruchrechnung-basis-k1-v1 (Original 2014-OS-B1d)
\begin{aufgabe}{Ein Spielwürfel wird einmal geworfen. Wie groß ist die Wahrscheinlichkeit für eine Zahl größer als 4? \hfill (P14) \\ \leerfeld}
\end{aufgabe}

\begleitteil
\erg{1}{P, denn $\frac{3}{8}$ – dort sind $3$ von $8$ gleichen Teilen grau}
\erg{2}{$\frac{1}{4} = 0{,}25$; die anderen sind $0{,}2$, $0{,}02$ und $0{,}04$}
\erg{3}{$2 \cdot (-1 - 7) = 2 \cdot (-8) = -16$}
\erg{4}{$0{,}3 \cdot 60 = 18\,€$}
\erg{5}{$x - 2 = 3$, also $x = 5$}
\erg{6}{B (eine Gerade ohne Knick)}
\erg{7}{2}
\erg{8}{$90 : 360 = 0{,}25 = 25\,\%$}
\erg{9}{$12$ Kanten ($4$ oben, $4$ unten, $4$ senkrecht)}
\erg{10}{$\frac{2}{6} = \frac{1}{3}$ (die 5 und die 6)}
\end{document}
